%% Resumo em Português
\begin{resumo}
O presente trabalho documenta o levantamento de requisitos, a modelagem arquitetural e a implementação da Prova de Conceito (PoC) do sistema \textbf{Karaokê Apache}, uma plataforma integrada para gerenciamento de fila, pedidos musicais e reprodução audiovisual desenvolvida para o estabelecimento comercial \textbf{Apache Gru Rock Bar} (\texttt{APACHE GRU ROCK BAR LTDA - ME}, CNPJ 64.687.141/0001-07), localizado em Guarulhos -- SP. O bar realiza semanalmente a \textit{Noite de Karaokê}, cujo controle de inscrições e escolha de faixas ocorria de forma manual, gerando perda de pedidos, duplicidade de músicas na fila, sobrecarga operacional do organizador ao alternar manualmente janelas de vídeo no telão e falta de visibilidade para os clientes sobre a ordem de apresentação. Como solução tecnológica, foi projetada uma arquitetura distribuída composta por quatro subsistemas desacoplados: (1) uma interface web responsiva para clientes (\textit{Cliente Web}), acessível via navegador móvel por QR Code na mesa para consulta do catálogo e envio de pedidos; (2) um painel administrativo (\textit{Admin Web}) para moderação, reordenação e controle da fila; (3) uma API Orquestradora centralizada responsável pela aplicação das regras de negócio (RN00 a RN04), autenticação e persistência no banco de dados relacional; e (4) o aplicativo desktop nativo \textbf{Karaokê TV Player}, construído em Electron e React com Vite, dedicado à exibição contínua no telão do bar. A metodologia de gestão combinou práticas ágeis (Scrum e Kanban via Trello e GitHub) com o desenvolvimento orientado a funcionalidades (\textit{Feature-Driven Development} -- FDD), estruturando os componentes \textit{Usuários}, \textit{Fila}, \textit{Pedido de Música} e \textit{Reprodutor}. Os testes da PoC demonstraram a sincronização em tempo real entre a interface de controle e o reprodutor Electron, validando a viabilidade técnica e econômica da solução para estabelecimentos de pequeno porte.

\vspace{\onelineskip}
\noindent
\textbf{Palavras-chaves}: Sistema web; Karaokê; Gerenciamento de fila; Electron; React; Projeto de extensão.
\end{resumo}
